\subsubsection{Vertical offset}
All results so far assume co-altitude flight. Figure~\ref{fig:X1} adds the vertical offset $d_z$: the benefit is concentrated in a narrow band around the wake sheet. At $d_y=0.85b$ the saving halves at $|d_z|\approx\ndzHalfFree\,b$ (\SI{\ndzHalfM}{m}) for the free wake and $\ndzHalfSheet\,b$ for the planar sheet; at $d_z=0.2b$ it is \ndzAtPtTwoFree\,\% (free) and \ndzAtPtTwoSheet\,\% (sheet) and at $0.4b$ only \ndzAtPtFourFree\,\% and \ndzAtPtFourSheet\,\%, against peaks of \ndzPeakFree\,\% and \ndzPeakSheet\,\%. The vertical station-keeping requirement, roughly $\pm0.1b$ ($\pm\SI{0.2}{m}$), is therefore as demanding as the lateral one, and altitude hold of small UAVs (barometer and GNSS noise of the order of decimetres) is a practical limit that the 2-D guidance model does not represent. Both models place the contour of zero benefit at $d_y\approx0.55$--$0.6b$ at $d_z=0$, and the zero contour widens with $|d_z|$: a follower that is far above or below the wake never gains, but it also never suffers the downwash penalty.
\begin{figure}[H]\centering\includegraphics[width=\textwidth]{X1_vertical_offset}
\caption{Follower electrical-power saving over lateral and vertical offset ($d_x=3b$, $\Va=\SI{18}{m/s}$) for (a) the planar sheet and (b) the free wake; (c) vertical cut at $d_y=0.85b$. Black: zero-saving contour; dot: nominal slot.}\label{fig:X1}\end{figure}

\subsubsection{Station-keeping accuracy}
The saving realised in flight is the expectation of the static saving over the slot error. Convolving the Gaussian lateral error distribution $\mathcal N(0,\sigma_y^2)$ with the static saving curve (evaluated on the Cauchy-midpoint lattice $d_y/b=m/60$ via 61-point Gauss--Hermite quadrature so collocation points never land on singular trailing filaments; Fig.~\ref{fig:X4}) shows how fast accuracy is lost: at $d_y=0.85b$ the VLM-sheet saving falls monotonically from \nskZeroSheet\,\% to \nskFifteenSheet\,\% at $\sigma_y=\SI{0.15}{m}$ (the NMPC's RMS error in turbulence), i.e.\ \nskRetain\,\% of the ideal value is retained, to \nskFortyKSheet\,\% at \SI{0.4}{m} and to \nskSixtySheet\,\% at \SI{0.6}{m}. The free-wake and strip models behave alike (\nskZeroFree\,\% $\to$ \nskFifteenFree\,\% $\to$ \nskFortyKFree\,\% and \nskZeroStrip\,\% $\to$ \nskFifteenStrip\,\% $\to$ \nskFortyKStrip\,\%). With GNSS-only relative navigation (error of order \SI{1}{m}) the expected benefit is negligible, which is the quantitative reason that the supervisor widens the lateral spacing when the relative-position uncertainty grows (Alg.~\ref{alg:sup}).
\begin{figure}[H]\centering\includegraphics[width=\textwidth]{X4_station_keeping_error}
\caption{Expected follower electrical-power saving versus lateral station-keeping error $\sigma_y$ (Gauss--Hermite quadrature over the Cauchy-midpoint static saving curve). (a) Three aerodynamic models at $d_y=0.85b$; (b) three nominal slots for the VLM planar sheet.}\label{fig:X4}\end{figure}

\subsubsection{Power curve and the choice of cruise speed}
Figure~\ref{fig:X6} places the wake benefit in the context of the aircraft's power curve. Electrical power is minimal at \SI{\npowMinV}{m/s} (\SI{\npowMinP}{W}) and energy per kilometre at \SI{\npowRangeV}{m/s} (\SI{\npowRangeE}{Wh/km}); the nominal \SI{18}{m/s} costs \SI{\npowEOneEight}{Wh/km}, only \npowExcess\,\% above the best-range value, so formation transit at a lower speed would not be competitive with simply cruising at the best-range speed unless the wake saving exceeds that margin. The percentage saving rises steeply as speed falls (induced-drag share grows), reaching \nsheetVThirteen\,\% (sheet) at \SI{13}{m/s}; in absolute terms, however, the power a follower can shed is at most about \SI{9}{W} at any speed, which is why the benefit cannot compete with the fixed overheads of reconfiguration unless the transit is long.
\begin{figure}[H]\centering\includegraphics[width=\textwidth]{X6_power_curve}
\caption{(a) Isolated electrical and aerodynamic power versus airspeed with the minimum-power point; (b) energy per kilometre and best-range speed; (c) follower saving at $d_y=0.85b$ versus airspeed for three wake models.}\label{fig:X6}\end{figure}
